\section{Multi-provider 支持与未来方向}

\subsection{Provider 无关的 LLM 抽象}

Utah 不直接调用 Anthropic SDK，而是使用 \texttt{pi-ai}——一个 provider-agnostic 的 LLM 抽象层，支持 Anthropic、OpenAI 和 Google。切换 provider 只需改配置：

\begin{lstlisting}[caption={LLM Provider 配置},language=Java]
llm: {
  provider: "anthropic", // or "openai" or "google"
  model: "claude-sonnet-4-20250514",
},
\end{lstlisting}

这种设计为未来的 Sub-agent 异构模型分配打下了基础——例如 coding Sub-agent 使用 Codex，research Sub-agent 使用 Opus。

\subsection{尚未解决的问题}

文章作者坦诚地列出了几个尚未完美解决的挑战：

\begin{table}[H]
\centering
\caption{Utah 当前的已知局限}
\begin{tabular}{lp{9cm}}
\toprule
\textbf{问题} & \textbf{现状} \\
\midrule
Steering（转向） & 新消息到达时 cancel+restart，进行中的工作丢失 \\
Streaming & 当前 step 是原子性的，不支持流式中间结果 \\
实时更新 & 尚未利用 Inngest 的 realtime 功能推送循环中间进度 \\
\bottomrule
\end{tabular}
\end{table}

\subsection{未来探索方向}

Utah 目前是一个运行在本地机器或服务器上的单人 harness。Inngest 团队正在探索将其扩展为\textbf{多人协作}的方向：

\begin{enumerate}
    \item \textbf{Swappable Sandboxes}：可替换的沙箱环境，使 Utah 能在 serverless 环境运行
    \item \textbf{External State and Memory}：外部化状态和记忆存储
    \item \textbf{Session Monitoring}：基于 Inngest API 和 Insights 功能的 coding session 监控
    \item \textbf{Human-in-the-Loop}：使用 \texttt{step.waitForEvent()} 创建人工审批流程
    \item \textbf{Self-building Agent}：让 Utah 能够编写自身——创建新 Agent、新 workflow、新 webhook，并通过 API 监控自身
\end{enumerate}

\begin{knowledgebox}{从 Single-player 到 Multi-player}
Utah 当前的架构——事件驱动、触发解耦、函数组合——天然支持向多人协作方向演进。多个 Agent 可以通过事件系统通信，每个 Agent 的 sandbox 可以独立替换，状态可以外部化到共享存储。这不需要架构上的根本改变，只需要在现有基础设施上\textbf{叠加}新能力。
\end{knowledgebox}

\subsection{本章小结}

Provider 无关的 LLM 抽象使得 Utah 能够灵活选择模型。文章坦诚地列出了当前的局限（steering、streaming、realtime），并描绘了从单人 harness 到多人协作平台的演进路线。这些方向进一步印证了 harness 思维的前瞻性——\textbf{好的基础设施设计天然为未来扩展留出空间}。


